\documentclass[9pt]{beamer}
\mode<presentation>

% Theme choice:
\usetheme{Madrid}%Darmstadt 

\usecolortheme[RGB={0,100,50}]{structure}
 \usefonttheme{structurebold} 
\setbeamercovered{invisible}
%\setbeamertemplate{navigation symbols}{}
\usepackage{dynblocks}
\usepackage{textpos}
\usepackage{amsfonts}
\usepackage{amsmath}
\usepackage{amssymb}
\usepackage{amsthm}

\newtheorem{teorema}{Teorema}\usepackage{mathtools}
\renewcommand{\proofname}{Prova}
\usepackage[portuguese]{babel}

\uselanguage{Portuguese}
\languagepath{Portuguese}

\deftranslation[to=Portuguese]{Theorem}{Teorema}
\deftranslation[to=Portuguese]{Lemma}{Lema}
\deftranslation[to=Portuguese]{Corollary}{Corolário}
\deftranslation[to=Portuguese]{Proof}{Prova}

\usepackage{dutchcal}
\usepackage{hyperref}
\usepackage[utf8]{inputenc}
\usepackage[T1]{fontenc}
\usepackage{amsthm}


\usepackage[affil-it]{authblk} 
\usepackage{etoolbox}
\usepackage{lmodern}

\makeatletter
\patchcmd{\@maketitle}{\LARGE \@title}{\fontsize{18}{19.2}\selectfont\@title}{}{}
\makeatother

\renewcommand\Authfont{\fontsize{12}{14.4}\selectfont}
\renewcommand\Affilfont{\fontsize{9}{10.8}\itshape}

% Title page details: 
\title[Presentation]{Lógica para Engenharia de Software} %add title


\institute[UniRV]{\textbf{\large  \textcolor{structure}{\textbf{\large
Análise Combinatória: Princípios de Contagem, Permutação,  Combinação e Arranjo.
}} %add name 
% \hspace{0.5pt}
\\ \qquad
\\Professor: Me. Sergio Souza Novak 
\\Faculdade de Engenharia de Software}
} 
% \titlegraphic{\includegraphics[width=5cm]{msu.eps}}
\titlegraphic{\includegraphics[width=5cm]{Universidade_de_Rio_Verde.jpg}}
\date{}

% Logo only on title page
%\logo{\includegraphics[width=1cm,keepaspectratio]{Logo.png}}
\usepackage[pscoord]{eso-pic}


\begin{document}

% Title page frame
%\begin{frame}
    %\titlepage 
%\end{frame}
\frame[plain]{\titlepage}
%logo
\AddToShipoutPictureFG{
    \put(\LenToUnit{.92\paperwidth},
    \LenToUnit{.185\paperheight})
    {\vtop{{\null}
    \makebox{\includegraphics[width=.8cm,keepaspectratio]{Universidade_de_Rio_Verde.jpg}}}}}


\begin{frame}{Sumário}
  \tableofcontents
\end{frame}

\section{Introdução}
\begin{frame}{Introdução}
Combinatória é o ramo da matemática que trata de contagem.
Esses problema são importantes quando temos recursos finitos, por exemplo,
\begin{itemize}
    \item Quanto espaço de armazenamento um determinado banco de dados usa?
    \item Quantos usuários uma determinada configuração de computador pode suportar?
    \item Quantos cálculos são efetuados por um determinado algoritmo?
    \item Quantas linhas tem uma tabela-verdade com $n$ letras de proposição?
\end{itemize}
\end{frame}

\section{Princípio da Multiplicação}
\begin{frame}{Princípio da Multiplicação}
\textbf{Exemplo:} uma criança pode escolher uma entre duas balas, uma rosa e outra preta, e um entre três chicletes, um amarelo, outro verde e outro branco. Quantos conjuntos diferentes a criança pode ter?

Para resolver este problema podemos fazer uma árvore de possibilidades de escolha.

\begin{center}
\begin{tikzpicture}[level distance=2.5cm,
  level 1/.style={sibling distance=3.5cm},
  level 2/.style={sibling distance=1cm}]
  \node {Início}
    child {node {Bala Rosa}
      child {node {Amarelo}}
      child {node {Verde}}
      child {node {Branco}}
    }
    child {node {Bala Preta}
      child {node {Amarelo}}
      child {node {Verde}}
      child {node {Branco}}
    };
\end{tikzpicture}
\end{center}

Sendo assim, a criança tem 6 possibilidades. A criança poderia escolher primeiro o chiclete e depois a bala, e o resultado seria o mesmo.
\end{frame}
\begin{frame}{Princípio Fundamental da Contagem (Princípio da Multiplicação)}

Sejam $E_1$ e $E_2$ dois experimentos finitos e não vazios, com conjuntos de resultados associados $A_1$ e $A_2$, tais que:
\[
|A_1| = n_1 \quad \text{e} \quad |A_2| = n_2
\]

\vspace{0.3cm}

Considere o experimento composto definido pelo produto cartesiano:
\[
A = A_1 \times A_2 = \{(a_1, a_2) \mid a_1 \in A_1 \land a_2 \in A_2\}
\]

\vspace{0.4cm}
\end{frame}
\begin{frame}{Princípio Fundamental da Contagem (Princípio da Multiplicação)}

\begin{block}{Formulação Lógica}
\[
\forall a_1 \in A_1, \; \exists n_2 \text{ possibilidades em } A_2
\;\;\Rightarrow\;\;
|A| = |A_1 \times A_2| = n_1 \cdot n_2
\]
\end{block}

\vspace{0.4cm}

De forma geral, para uma sequência finita de conjuntos $A_1, A_2, \dots, A_k$, com $|A_i| = n_i$, tem-se:
\[
\left| \prod_{i=1}^{k} A_i \right| = \prod_{i=1}^{k} n_i
\]

\vspace{0.2cm}

onde:
\[
\prod_{i=1}^{k} A_i = \{(a_1, \dots, a_k) \mid \forall i \in \{1,\dots,k\}, \; a_i \in A_i\}
\]

\end{frame}
\begin{frame}{Princípio Fundamental da Contagem}

\textbf{Exercício 1:} Considere a formação de sequências de 4 dígitos utilizando o conjunto $D = \{0,1,2,3,4,5,6,7,8,9\}$. Admitindo-se a repetição de dígitos, quantas sequências distintas de comprimento 4 podem ser formadas?

\vspace{0.5cm}

\textbf{Exercício 2:} Nas mesmas condições do exercício anterior, determine o número de sequências de 4 dígitos que podem ser formadas sob a restrição de que não haja repetição de dígitos.

\end{frame}
\begin{frame}{Resolução com Formalismo Matemático}

\textbf{Exercício 1:}

Seja $D = \{0,1,2,3,4,5,6,7,8,9\}$ o conjunto dos dígitos possíveis, com $|D| = 10$.

Definimos o conjunto de todos os números de 4 dígitos (permitindo repetição) como:
\[
A = D \times D \times D \times D
\]

Logo,
\[
|A| = |D|^4 = 10^4 = 10000
\]

\vspace{0.4cm}

\textbf{Exercício 2:}

Agora, consideramos que não há repetição de dígitos. Assim, definimos:
\[
A = \{(d_1,d_2,d_3,d_4) \in D^4 \mid \forall i,j \in \{1,2,3,4\},\; i \neq j \Rightarrow d_i \neq d_j\}
\]

\vspace{0.3cm}

Pelo Princípio da Multiplicação:
\[
|A| = 10 \cdot 9 \cdot 8 \cdot 7 = 5040
\]


\end{frame}

\section{Princípio da Adição}
\begin{frame}{Princípio da Adição}
\textbf{Exemplo:} suponha que queremos selecionar uma sobremesa entre três tortas e quatro bolos. De quantas maneiras podemos fazer isso?
\vspace{0.3cm}

Temos dois eventos, um com 3 resultados e outro com 4 resultados. No entanto, nesse caso não podemos multiplicar, pois queremos comer apenas uma sobremesa.
\vspace{0.1cm}

\begin{center}
\begin{tikzpicture}
  \node[draw, rounded corners, minimum width=3cm, minimum height=1.5cm] (tortas) at (-2,0) {3 Tortas};
  \node[draw, rounded corners, minimum width=3cm, minimum height=1.5cm] (bolos) at (2,0) {4 Bolos};
  \node at (0, 0) {\textbf{OU}};
\end{tikzpicture}
\end{center}

O número de escolhas possíveis é o número total de escolhas que temos: $3+4=7$.
\end{frame}
\begin{frame}{Princípio Fundamental da Contagem (Princípio da Adição)}

Sejam $A$ e $B$ conjuntos finitos que representam espaços de resultados de dois eventos.

\vspace{0.3cm}

\begin{block}{Definição Formal}
Se $A$ e $B$ são disjuntos, isto é,
\[
A \cap B = \emptyset
\]
então:
\[
|A \cup B| = |A| + |B|
\]
\end{block}


\begin{block}{Formulação Lógica}
\[
\forall x \; (x \in A \Rightarrow x \notin B)
\;\;\Rightarrow\;\;
|A \cup B| = |A| + |B|
\]
\end{block}


De forma geral, para uma família finita de conjuntos disjuntos $\{A_i\}_{i=1}^k$, tal que:
\[
\forall i \neq j,\; A_i \cap A_j = \emptyset
\]
tem-se:
\[
\left| \bigcup_{i=1}^{k} A_i \right| = \sum_{i=1}^{k} |A_i|
\]

\end{frame}

\begin{frame}{Princípio da Adição}
\textbf{Exercício 3:} Um consumidor deseja comprar um veículo de uma concessionária. A concessionária tem 23 automóveis e 14 caminhões em estoque. Quantas escolhas possíveis o consumidor tem?

\vspace{0.5cm}

\textbf{Obs:} poderíamos utilizar o princípio da adição se o consumidor quisesse comprar um veículo de uma concessionária com 23 automóveis, 14 caminhões e 17 veículos vermelhos?

\vspace{0.5cm}
Lembre que: Se $A$ e $B$ são conjuntos finitos disjuntos, então $|A \cup B| = |A| + |B|$.
\end{frame}

\begin{frame}{Princípio da Adição}
\textbf{Exercício 4:} Na linguagem de programação BASIC original, um identificador tem que ser uma letra ou uma letra seguida de um único dígito. Quantos identificadores podemos formar?
\end{frame}

\begin{frame}{Resolução}
\textbf{Exercício 3:}
Seja $A$ o conjunto de automóveis e $C$ o conjunto de caminhões.
Temos $|A| = 23$ e $|C| = 14$. Como não é possível um veículo ser carro e caminhão ao mesmo tempo, $A \cap C = \emptyset$.
Pelo Princípio da Adição:
\[ |A \cup C| = |A| + |C| = 23 + 14 = 37 \text{ escolhas possíveis.} \]

\textbf{Obs:} Se houvesse 17 veículos vermelhos (conjunto $V$), não poderíamos simplesmente somar, pois os conjuntos não seriam disjuntos (há carros e caminhões vermelhos). Teríamos que usar o Princípio da Inclusão-Exclusão: $|A \cup C \cup V|$.

\vspace{0.4cm}

\textbf{Exercício 4:}
Um identificador pode ter dois formatos (disjuntos):
\begin{itemize}
    \item \textbf{Formato 1:} Apenas uma letra. Quantidade = 26.
    \item \textbf{Formato 2:} Uma letra e um dígito. Pelo Princípio da Multiplicação: $26 \cdot 10 = 260$.
\end{itemize}
Pelo Princípio da Adição, o total é $26 + 260 = 286$ identificadores.
\end{frame}

\section{Princípio da Adição e Multiplicação}
\begin{frame}{Princípio da Adição e Multiplicação}
Como podemos ver no exercício 4, o princípio da adição é usado junto com o princípio da multiplicação.
\vspace{0.3cm}

Considerando o exemplo das ``balas e chicletes'', escolher uma bala rosa e depois um chiclete amarelo não é a mesma coisa que escolher um chiclete amarelo e depois uma bala rosa.
\vspace{0.3cm}

Podemos considerar dois casos disjuntos: cada um desses casos (princípio da multiplicação) tem 6 possibilidades e pelo princípio da adição existem $6+6=12$ maneiras diferentes de escolher os doces.
\end{frame}

\begin{frame}{Princípio da Adição e Multiplicação}
\textbf{Exemplo 5:} quantos números de quatro dígitos começam com 4 ou 5?

\vspace{0.4cm}

\textbf{Exemplo 6:} quantos inteiros de três dígitos são pares?

\vspace{0.4cm}

\textbf{Exemplo 7:} suponha que os quatro últimos dígitos de um número de telefone têm que incluir pelo menos um dígito repetido. Quantos desses números existem? (Dica: pensar em conjuntos disjuntos)
\end{frame}

\begin{frame}{Resolução}
\textbf{Exemplo 5:}
Seja $A$ o conjunto dos que começam com 4, e $B$ os que começam com 5. O primeiro dígito é fixo.
$|A| = 1 \cdot 10 \cdot 10 \cdot 10 = 1000$ e $|B| = 1 \cdot 10 \cdot 10 \cdot 10 = 1000$.
Como $A \cap B = \emptyset$, o total é $1000 + 1000 = 2000$.

\vspace{0.3cm}

\textbf{Exemplo 6:}
O primeiro dígito não pode ser 0 (9 opções: 1 a 9). O segundo dígito pode ser qualquer um (10 opções). O último dígito deve ser par, ou seja, $\in \{0,2,4,6,8\}$ (5 opções).
Total: $9 \cdot 10 \cdot 5 = 450$.

\vspace{0.3cm}

\textbf{Exemplo 7:}
É mais fácil contar pelo complemento.
Total de combinações de 4 dígitos: $10^4 = 10000$.
Total de combinações de 4 dígitos \textbf{sem repetição} (visto antes): $10 \cdot 9 \cdot 8 \cdot 7 = 5040$.
Pelo princípio da diferença, as combinações que têm \textbf{pelo menos uma} repetição são:
$10000 - 5040 = 4960$.
\end{frame}

\section{Permutação}
\begin{frame}{Permutação}
Lembre o exercício 2, sobre os quatro últimos dígitos de um número de telefone não repetidos. Perceba que o número 1259 não é igual ao número 2951, já que a ordem dos dígitos é importante.
\vspace{0.3cm}

Um arranjo ordenado de objetos é chamado de \textbf{permutação}.
\vspace{0.1cm}

\begin{center}
\begin{tikzpicture}
  \node[draw, minimum width=0.8cm, minimum height=0.8cm] (a1) at (0,0) {1};
  \node[draw, minimum width=0.8cm, minimum height=0.8cm] (a2) at (1,0) {2};
  \node[draw, minimum width=0.8cm, minimum height=0.8cm] (a3) at (2,0) {5};
  \node[draw, minimum width=0.8cm, minimum height=0.8cm] (a4) at (3,0) {9};
  \node at (4.5, 0) {$\neq$};
  \node[draw, minimum width=0.8cm, minimum height=0.8cm] (b1) at (6,0) {2};
  \node[draw, minimum width=0.8cm, minimum height=0.8cm] (b2) at (7,0) {9};
  \node[draw, minimum width=0.8cm, minimum height=0.8cm] (b3) at (8,0) {5};
  \node[draw, minimum width=0.8cm, minimum height=0.8cm] (b4) at (9,0) {1};
\end{tikzpicture}
\end{center}

O número de permutações de $r$ objetos distintos escolhidos entre $n$ objetos distintos é denotado $P(n,r)$.
\vspace{0.1cm}

Portanto, a solução do problema dos 4 dígitos sem repetição pode ser expressada por $P(10, 4)$.
\end{frame}


\begin{frame}{Permutações (Formalismo)}

Seja $A$ um conjunto finito tal que $|A| = n$.

\vspace{0.3cm}

\begin{block}{Definição}
Uma \textbf{permutação} de $A$ é uma bijeção:
\[
f: A \rightarrow A
\]
ou, equivalentemente, uma sequência ordenada contendo todos os elementos de $A$ sem repetição.
\end{block}

\vspace{0.4cm}

O número total de permutações de $A$ é dado por:
\[
P(n) = n!
\]


\end{frame}


\begin{frame}{Permutações (Aplicações)}

\textbf{Exemplo 1:} Seja $A = \{a,b,c\}$. O número de permutações é:
\[
P(3) = 3! = 6
\]

\vspace{0.3cm}

Explicitamente:
\[
\{(a,b,c), (a,c,b), (b,a,c), (b,c,a), (c,a,b), (c,b,a)\}
\]

\vspace{0.5cm}

\textbf{Exemplo 2:} Seja $A = \{\text{c, o, m, p, i, l, a, r}\}$, com $|A| = 8$.

\vspace{0.3cm}

Número de sequências ordenadas de 3 elementos distintos:
\[
P(8,3) = 8 \cdot 7 \cdot 6 = 336
\]

\vspace{0.4cm}

\textbf{Interpretação:}
\[
\forall a_1 \in A,\;
\forall a_2 \in A \setminus \{a_1\},\;
\forall a_3 \in A \setminus \{a_1,a_2\}
\]

\vspace{0.2cm}

\textbf{Observação:} A ordem dos elementos é relevante.

\end{frame}


\begin{frame}{Permutação}
\textbf{Exercício 1:} Calcule o valor de $P(n,0)$, $P(n,1)$ e $P(n,n)$.

\vspace{0.2cm}

\textbf{Exercício 2:} dez atletas competem num evento olímpico. São dadas medalhas de ouro, prata e bronze. De quantas maneiras diferentes podem ser dadas as medalhas?

\vspace{0.2cm}

\textbf{Exercício 3:} de quantas maneiras diferentes se pode selecionar um presidente e um vice-presidente de um grupo de 20 pessoas?

\vspace{0.2cm}

\textbf{Exercício 4:} uma biblioteca tem 4 livros sobre sistemas operacionais, 7 sobre programação e 3 sobre estrutura de dados. De quantas maneiras diferentes esses livros podem ser arrumados em uma prateleira, dado que todos os livros sobre o mesmo assunto devem ficar juntos?
\end{frame}

\begin{frame}{Resolução}
\textbf{Exercício 1:}
\begin{itemize}
    \item $P(n, 0) = \frac{n!}{(n-0)!} = \frac{n!}{n!} = 1$
    \item $P(n, 1) = \frac{n!}{(n-1)!} = \frac{n \cdot (n-1)!}{(n-1)!} = n$
    \item $P(n, n) = \frac{n!}{(n-n)!} = \frac{n!}{0!} = \frac{n!}{1} = n!$
\end{itemize}

\vspace{0.1cm}
\textbf{Exercício 2:}
Importa a ordem (ouro $\neq$ prata). $P(10, 3) = 10 \cdot 9 \cdot 8 = 720$ maneiras.

\vspace{0.1cm}
\textbf{Exercício 3:}
Novamente a ordem importa (presidente $\neq$ vice). $P(20, 2) = 20 \cdot 19 = 380$ maneiras.

\vspace{0.1cm}
\textbf{Exercício 4:}
Primeiro, permutamos os $3$ assuntos entre si: $P(3,3) = 3! = 6$.
Dentro de cada assunto, permutamos os livros:
\begin{itemize}
    \item Sistemas Operacionais: $4! = 24$
    \item Programação: $7! = 5040$
    \item Estrutura de Dados: $3! = 6$
\end{itemize}
Total = $3! \cdot 4! \cdot 7! \cdot 3! = 6 \cdot 24 \cdot 5040 \cdot 6 = 4.354.560$ maneiras.
\end{frame}

\section{Arranjo}
\begin{frame}{Arranjo Simples}

\begin{block}{Definição}
Um \textbf{arranjo simples} de $n$ elementos tomados $p$ a $p$ (com $p \leq n$) é qualquer sequência ordenada de $p$ elementos distintos escolhidos de um conjunto de $n$ elementos.
\end{block}

\vspace{0.3cm}
O número total de arranjos é denotado por $A(n,p)$ ou $A_{n,p}$, e é calculado por:
\[
A(n, p) = \frac{n!}{(n-p)!}
\]

\vspace{0.2cm}
\textbf{Observação:} Diferente da combinação, no arranjo a \textbf{ordem importa}. Diferente da permutação simples vista na seção anterior ($P(n)=n!$), o arranjo em geral escolhe um subconjunto $p \leq n$.
\end{frame}

\begin{frame}{Arranjo (Diagrama)}
Exemplo para $n=4$ (conjunto $\{1,2,3,4\}$) e agrupamentos de $p=2$:
\vspace{0.2cm}
\begin{center}
\begin{tikzpicture}[
    scale=0.9,
    every node/.style={circle, draw, minimum size=0.8cm, inner sep=1pt}
]

% Conjunto n=4
\node (x1) at (0,0.5) {1};
\node (x2) at (1,0.5) {2};
\node (x3) at (0,-0.5) {3};
\node (x4) at (1,-0.5) {4};
\node[draw=none, rectangle] at (0.5,-1.3) {$n=4$ elementos};

\draw[->, thick] (2,0) -- (3.5,0) node[midway, above] {$p=2$};

% Arranjos
\node[rectangle, draw=none] (a1) at (5, 1.5) {$(1, 2) \neq (2, 1)$};
\node[rectangle, draw=none] (a2) at (5, 0.5) {$(1, 3) \neq (3, 1)$};
\node[rectangle, draw=none] (a3) at (5, -0.5) {$(1, 4) \neq (4, 1)$};
\node[rectangle, draw=none] (a4) at (8, 1.5) {$(2, 3) \neq (3, 2)$};
\node[rectangle, draw=none] (a5) at (8, 0.5) {$(2, 4) \neq (4, 2)$};
\node[rectangle, draw=none] (a6) at (8, -0.5) {$(3, 4) \neq (4, 3)$};

\node[draw=none, rectangle] at (6.5,-1.5) {$A(4,2) = 12$ sequências ordenadas};

\end{tikzpicture}
\end{center}
\end{frame}

\begin{frame}{Exercícios - Arranjo}
\textbf{Exercício 1:} Quantos números de 3 algarismos distintos podemos formar com os algarismos $\{1, 2, 3, 4, 5\}$?

\vspace{0.4cm}

\textbf{Exercício 2:} Em uma corrida com 10 carros, de quantas maneiras pode ser formado o pódio (1º, 2º e 3º lugares)?

\vspace{0.4cm}

\textbf{Exercício 3:} Uma senha de banco é formada por 4 caracteres distintos escolhidos entre as 26 letras do alfabeto. Quantas senhas possíveis existem?
\end{frame}

\begin{frame}{Resolução - Arranjo}
\textbf{Exercício 1:}
Temos $n = 5$ algarismos e queremos formar números com $p = 3$ algarismos distintos. A ordem importa ($123 \neq 321$).
\[
A(5,3) = \frac{5!}{(5-3)!} = \frac{5!}{2!} = \frac{5 \cdot 4 \cdot 3 \cdot 2!}{2!} = 60
\]
\vspace{0.2cm}

\textbf{Exercício 2:}
Temos $n = 10$ carros e queremos escolher $p = 3$ para o pódio. A ordem importa (1º $\neq$ 2º $\neq$ 3º).
\[
A(10,3) = \frac{10!}{(10-3)!} = \frac{10!}{7!} = 10 \cdot 9 \cdot 8 = 720
\]
\vspace{0.2cm}

\textbf{Exercício 3:}
Temos $n = 26$ letras e queremos uma senha de $p = 4$ caracteres distintos. A ordem importa.
\[
A(26,4) = \frac{26!}{(26-4)!} = \frac{26!}{22!} = 26 \cdot 25 \cdot 24 \cdot 23 = 358.800
\]
\end{frame}

\section{Combinação}
\begin{frame}{Combinação}

Algumas vezes queremos selecionar $r$ objetos de um conjunto de $n$ objetos, sem considerar a ordem. Nesse caso, estamos interessados no número de combinações, denotado por $\binom{n}{r}$.

\vspace{0.3cm}

Para cada subconjunto de $r$ elementos, existem $r!$ formas de ordená-los. Logo:
\[
\binom{n}{r} \cdot r! = P(n, r)
\]



\begin{center}
\begin{tikzpicture}[
    scale=0.9,
    every node/.style={circle, draw, minimum size=0.8cm}
]

% Combinação (sem ordem)
\node (a1) at (0,0) {A};
\node (a2) at (1.2,0) {B};
\node (a3) at (2.4,0) {C};

\node[draw=none, rectangle] at (1.2,-1) {Subconjunto $\{A,B,C\}$};

% Permutações (com ordem)
\node (p1) at (5,0.8) {ABC};
\node (p2) at (6.5,0.8) {ACB};
\node (p3) at (5,-0.2) {BAC};
\node (p4) at (6.5,-0.2) {BCA};
\node (p5) at (5,-1.2) {CAB};
\node (p6) at (6.5,-1.2) {CBA};

\node[draw=none, rectangle] at (5.75,-2) {$3! = 6$ ordenações};

\draw[->, thick] (2.8,0) -- (4.2,0) node[midway, above] {$r!$};

\end{tikzpicture}
\end{center}

\end{frame}
\begin{frame}{Combinação}
Ou seja,
\[
\binom{n}{r} = \frac{n!}{r!(n - r)!} \quad \text{para } 0 \leq r \leq n
\]

Outras notações: $C_{n,r}$, $C_n^r$ ou $\binom{n}{r}$.

\vspace{0.5cm}

\textbf{Exemplo:} Dez atletas competem em um evento olímpico; três deles serão declarados vencedores. De quantas maneiras podem ser escolhidos os vencedores?
\end{frame}

\begin{frame}{Combinação}
\textbf{Exercício 1:} Calcule o valor de $\binom{7}{3}$.

\vspace{0.3cm}

\textbf{Exercício 2:} Calcule o valor de $\binom{n}{0}$, $\binom{n}{1}$ e $\binom{n}{n}$.

\vspace{0.3cm}

\textbf{Exercício 3:} Uma comissão de 8 alunos deve ser escolhida em um grupo contendo 19 alunos do primeiro ano e 34 do segundo ano.
\begin{itemize}
    \item[a)] De quantas maneiras é possível selecionar 3 alunos do primeiro ano e 5 do segundo?
\end{itemize}
\end{frame}

\begin{frame}{Combinação}
\textbf{Exercício 3:} Uma comissão de 8 alunos deve ser escolhida em um grupo contendo 19 alunos do primeiro ano e 34 do segundo ano.
\begin{itemize}
    \item[b)] De quantas maneiras é possível selecionar uma comissão contendo exatamente 1 aluno do primeiro ano?
    \item[c)] De quantas maneiras é possível selecionar uma comissão contendo no máximo 1 aluno do primeiro ano?
    \item[d)] De quantas maneiras é possível selecionar uma comissão contendo pelo menos 1 aluno do primeiro ano?
\end{itemize}
\end{frame}

\begin{frame}{Resolução}
\textbf{Exemplo e Exercício 1, 2 e 3:}
\textbf{Exemplo atletas:} $\binom{10}{3} = \frac{10!}{3!7!} = \frac{10 \cdot 9 \cdot 8}{3 \cdot 2 \cdot 1} = 120$.

\vspace{0.2cm}
\textbf{Ex1:} $\binom{7}{3} = \frac{7!}{3!4!} = \frac{7 \cdot 6 \cdot 5}{3 \cdot 2 \cdot 1} = 35$.

\vspace{0.2cm}
\textbf{Ex2:}
\begin{itemize}
    \item $\binom{n}{0} = \frac{n!}{0!(n-0)!} = \frac{n!}{1 \cdot n!} = 1$
    \item $\binom{n}{1} = \frac{n!}{1!(n-1)!} = \frac{n \cdot (n-1)!}{(n-1)!} = n$
    \item $\binom{n}{n} = \frac{n!}{n!(n-n)!} = \frac{n!}{n!0!} = 1$
\end{itemize}

\vspace{0.2cm}
\textbf{Ex3:} $\binom{52}{5} = \frac{52!}{5!47!} = \frac{52 \cdot 51 \cdot 50 \cdot 49 \cdot 48}{120} = 2.598.960$.
\end{frame}

\begin{frame}{Resolução - Exercício 4}
Temos 19 alunos do 1º ano e 34 do 2º ano. Comissão de 8.

\vspace{0.2cm}

\textbf{a) 3 do 1º ano e 5 do 2º:}
Escolhemos 3 dos 19: $\binom{19}{3} = 969$.
Escolhemos 5 dos 34: $\binom{34}{5} = 278.256$.
Total = $\binom{19}{3} \cdot \binom{34}{5} = 969 \cdot 278.256 = 269.629.584$.

\vspace{0.2cm}

\textbf{b) Exatamente 1 do 1º ano:}
Escolhemos 1 dos 19: $\binom{19}{1} = 19$.
Escolhemos 7 dos 34: $\binom{34}{7} = 5.379.616$.
Total = $19 \cdot 5.379.616 = 102.212.704$.
\end{frame}

\begin{frame}{Resolução - Exercício 4 (Cont.)}
\textbf{c) No máximo 1 aluno do 1º ano:}
Isso significa "0 alunos do 1º ano" OU "1 aluno do 1º ano".
- Nenhum do 1º ano (todos do 2º): $\binom{19}{0} \cdot \binom{34}{8} = 1 \cdot 18.156.204 = 18.156.204$.
- 1 aluno do 1º ano (calculado na letra b): $102.212.704$.
Total = $18.156.204 + 102.212.704 = 120.368.908$.

\vspace{0.2cm}

\textbf{d) Pelo menos 1 aluno do 1º ano:}
Pelo complementar, é igual ao total de comissões menos as comissões onde \textbf{nenhum} aluno é do 1º ano.
Total de comissões possíveis (de 53 alunos, escolhemos 8):
$\binom{53}{8}$
Comissões com nenhum aluno do 1º ano (visto na letra c):
$\binom{34}{8}$
Total = $\binom{53}{8} - \binom{34}{8}$.
\end{frame}

\section{Eliminação de Duplicatas}
\begin{frame}{Eliminação de Duplicatas}
Em geral, suponha que temos $n$ objetos dos quais um conjunto de $n_1$ são indistinguíveis (idênticos) entre si, um outro conjunto de $n_2$ são indistinguíveis entre si, e assim por diante, até um conjunto de $n_k$ objetos indistinguíveis entre si. O número de permutações distintas desses $n$ objetos é:
\[
\frac{n!}{n_1! n_2! \dots n_k!}
\]
\vspace{0.3cm}

Pense na permutação da palavra \textbf{ARARA}.

\vspace{0.3cm}

\textbf{Exemplo:} Quantas permutações distintas podem ser feitas com os caracteres que formam a palavra \textbf{MISSISSIPPI}?
\end{frame}

\begin{frame}{Resolução}
\textbf{Palavra ARARA:}
Temos $n = 5$ letras no total.
Letra A repete $3$ vezes ($n_1 = 3$).
Letra R repete $2$ vezes ($n_2 = 2$).
Formas distintas: $\frac{5!}{3!2!} = \frac{120}{6 \cdot 2} = 10$.

\vspace{0.4cm}

\textbf{Palavra MISSISSIPPI:}
Temos $n = 11$ letras no total.
Frequências:
\begin{itemize}
    \item M: 1 vez
    \item I: 4 vezes
    \item S: 4 vezes
    \item P: 2 vezes
\end{itemize}
Formas distintas: $\frac{11!}{1! 4! 4! 2!} = \frac{39.916.800}{1 \cdot 24 \cdot 24 \cdot 2} = 34.650$.
\end{frame}
\section{Classificação dos Problemas}
\begin{frame}{Classificação dos Problemas de Contagem}
Ao resolver problemas de contagem, devemos responder duas perguntas:

\begin{itemize}
    \item A ordem importa?
    \item É permitido repetir elementos?
\end{itemize}

\begin{center}
\begin{tikzpicture}[
    level distance=1.5cm,
    level 1/.style={sibling distance=5.5cm},
    level 2/.style={sibling distance=2.5cm},
    every node/.style={rectangle, rounded corners, draw, align=center, fill=blue!10, font=\footnotesize},
    edge from parent/.style={draw, -latex}
]

\node[fill=blue!30, font=\bfseries] {A ordem importa?}
    child {node {Sim}
        child {node {Sem Repetição\\ \textbf{Arranjo}}}
        child {node {Com Repetição\\ \textbf{Arranjo com Repet.}}}
    }
    child {node {Não}
        child {node {Sem Repetição\\ \textbf{Combinação}}}
        child {node {Com Repetição\\ \textbf{Combinação com Repet.}}}
    };
\end{tikzpicture}
\end{center}
\end{frame}
\section{Arranjo com Repetição}
\begin{frame}{Arranjo com Repetição}

\begin{block}{Definição}
Um arranjo com repetição de $n$ elementos tomados $p$ a $p$ permite que um mesmo elemento seja escolhido mais de uma vez.
\end{block}

\vspace{0.3cm}

O número de arranjos é dado por:
\[
A'(n,p) = n^p
\]

\vspace{0.3cm}

\textbf{Exemplo:} Senhas de 4 dígitos (0–9).
\vspace{0.2cm}

\begin{center}
\begin{tikzpicture}
  \node[draw, minimum width=0.8cm, minimum height=0.8cm, fill=gray!20] (d1) at (0,0) {0-9};
  \node[draw, minimum width=0.8cm, minimum height=0.8cm, fill=gray!20] (d2) at (1.5,0) {0-9};
  \node[draw, minimum width=0.8cm, minimum height=0.8cm, fill=gray!20] (d3) at (3,0) {0-9};
  \node[draw, minimum width=0.8cm, minimum height=0.8cm, fill=gray!20] (d4) at (4.5,0) {0-9};
  
  \node[draw=none, rectangle] at (0,-0.8) {10};
  \node[draw=none, rectangle] at (0.75,-0.8) {$\times$};
  \node[draw=none, rectangle] at (1.5,-0.8) {10};
  \node[draw=none, rectangle] at (2.25,-0.8) {$\times$};
  \node[draw=none, rectangle] at (3,-0.8) {10};
  \node[draw=none, rectangle] at (3.75,-0.8) {$\times$};
  \node[draw=none, rectangle] at (4.5,-0.8) {10};
  
  \node[draw=none, rectangle] at (2.25,-1.5) {$10^4 = 10000$ senhas};
\end{tikzpicture}
\end{center}

\end{frame}
\section{Combinação com Repetição}
\begin{frame}{Combinação com Repetição}

\begin{block}{Definição}
Na combinação com repetição, podemos escolher elementos repetidos, e a ordem não importa.
\end{block}

\vspace{0.3cm}

O número de combinações é dado por:
\[
C'(n,p) = \binom{n + p - 1}{p}
\]

\vspace{0.1cm}

\textbf{Exemplo:} Escolher 3 sorvetes entre 5 sabores (podendo repetir):
\begin{center}
\begin{tikzpicture}[scale=0.9]
  % Exemplo
  \node[circle, fill=blue!50, inner sep=2.5pt] at (0,0) {};
  \node[draw=none] at (0.5,0) {\Large $|$};
  \node[draw=none] at (1,0) {\Large $|$};
  \node[circle, fill=blue!50, inner sep=2.5pt] at (1.5,0) {};
  \node[circle, fill=blue!50, inner sep=2.5pt] at (2,0) {};
  \node[draw=none] at (2.5,0) {\Large $|$};
  \node[draw=none] at (3,0) {\Large $|$};
  
  \node[draw=none, font=\footnotesize] at (0,-0.6) {S1};
  \node[draw=none, font=\footnotesize] at (0.75,-0.6) {S2};
  \node[draw=none, font=\footnotesize] at (1.75,-0.6) {S3};
  \node[draw=none, font=\footnotesize] at (2.75,-0.6) {S4};
  \node[draw=none, font=\footnotesize] at (3.5,-0.6) {S5};
  
  \node[draw=none, text width=5cm, align=center, font=\footnotesize] at (1.5, -1.5) {3 Escolhas ($\bullet$) + 4 Separadores ($|$)\\ = 7 posições em fila};
  
  \node[draw=none, font=\small] at (6, -0.6) {$C'(5,3) = \binom{3+5-1}{3} = \binom{7}{3} = 35$};
\end{tikzpicture}
\end{center}

\end{frame}
\begin{frame}{Intuição: Com ou Sem Repetição}

\begin{center}
\begin{tikzpicture}
  \node[draw, rectangle, rounded corners, fill=red!10, text width=4.5cm, align=center, minimum height=2.5cm] (sem) at (0,0) {
    \textbf{Sem Repetição} \\ \vspace{0.2cm}
    Um elemento escolhido \textbf{não volta} para o conjunto. \\ \vspace{0.2cm}
    \textit{Ex: Equipes, Pódios, Senhas c/ letras distintas}
  };
  
  \node[draw, rectangle, rounded corners, fill=green!10, text width=4.5cm, align=center, minimum height=2.5cm] (com) at (6,0) {
    \textbf{Com Repetição} \\ \vspace{0.2cm}
    Um elemento escolhido \textbf{continua disponível}. \\ \vspace{0.2cm}
    \textit{Ex: Senhas numéricas comuns, Sabores de sorvete}
  };
\end{tikzpicture}
\end{center}

\vspace{0.3cm}

\textbf{Dica prática:}
Se o problema não proíbe repetição, geralmente ela é permitida.
\end{frame}
\end{document}

